\documentclass[conference]{IEEEtran}
\IEEEoverridecommandlockouts
\usepackage{cite}
\usepackage{amsmath,amssymb,amsfonts}
\usepackage{algorithmic}
\usepackage{graphicx}
\usepackage{textcomp}
\usepackage{xcolor}
\usepackage{booktabs}
\usepackage{multirow}
\usepackage{array}
\usepackage{url}
\def\BibTeX{{\rm B\kern-.05em{\sc i\kern-.025em b}\kern-.08em
    T\kern-.1667em\lower.7ex\hbox{E}\kern-.125emX}}

\begin{document}

\title{B-CIRL-TFM: Concurrent Imitation and Reinforcement Learning\\
for Automated Data Cleaning with Calibrated Predictions}

\author{
\IEEEauthorblockN{Aryan Nangarath}
\IEEEauthorblockA{\textit{Dept. of Computer Science} \\
\textit{PES University}\\
Bengaluru, India}
\and
\IEEEauthorblockN{Dhanush Shekhar}
\IEEEauthorblockA{\textit{Dept. of Computer Science} \\
\textit{PES University}\\
Bengaluru, India}
\and
\IEEEauthorblockN{Abhishek Aithal}
\IEEEauthorblockA{\textit{Dept. of Computer Science} \\
\textit{PES University}\\
Bengaluru, India}
}

\maketitle

\begin{abstract}
Automated data cleaning remains a critical bottleneck in machine learning
pipelines. While reinforcement learning (RL) approaches such as Learn2Clean
have demonstrated the ability to discover optimal cleaning sequences, they
suffer from slow convergence due to cold-start random initialization and
produce poorly calibrated predictions. We propose \textbf{B-CIRL-TFM}
(Behavioural Concurrent Imitation-Reinforcement Learning with TabPFN),
a novel framework that injects a decaying Behavioural Cloning (BC) loss
alongside standard Proximal Policy Optimization (PPO) updates. This
concurrent training prevents catastrophic forgetting of expert knowledge
while allowing the agent to discover dataset-specific improvements through
reward signals. We further introduce an ECE proxy penalty to regularize
calibration during cleaning, and a binary action gating mechanism that
restricts harmful operations on categorical data. Evaluated on 12 benchmark
datasets, B-CIRL-TFM achieves a mean ECE of 0.0512, a 17\% improvement
over the uncleaned baseline (0.0618), and outperforms all static baselines
on model calibration while matching or exceeding accuracy on 8 of 12
datasets.
\end{abstract}

\begin{IEEEkeywords}
data cleaning, imitation learning, reinforcement learning, behavioural cloning,
TabPFN, calibration, ECE, automated machine learning
\end{IEEEkeywords}

\section{Introduction}

Data preparation is widely recognized as consuming 60--80\% of a data
scientist's time~\cite{dasu2003exploratory}. Despite this, most automated
machine learning (AutoML) pipelines focus on model selection and
hyperparameter tuning, leaving data cleaning as a manual, ad hoc process.

Learn2Clean~\cite{berti2019learn2clean} introduced the paradigm of
formulating data cleaning as a sequential decision-making problem solvable
by reinforcement learning (RL). The agent learns to select, order, and
parameterize cleaning operations to maximize downstream model performance.
However, RL-based cleaning suffers from two fundamental limitations:

\begin{enumerate}
    \item \textbf{Cold-start problem}: The agent begins with a random policy,
    requiring thousands of environment interactions before learning meaningful
    cleaning strategies. This is particularly wasteful because domain
    experts already know that imputation should precede outlier removal,
    which should precede scaling.
    \item \textbf{Calibration neglect}: Existing approaches optimize for
    accuracy or F1, ignoring Expected Calibration Error (ECE). Poorly
    calibrated models produce overconfident predictions that are harmful
    in high-stakes domains such as medical diagnosis.
\end{enumerate}

We address both limitations through \textbf{B-CIRL-TFM}, which makes the
following contributions:

\begin{itemize}
    \item A \textit{concurrent} IL+RL training architecture that injects
    a linearly decaying BC loss at every PPO gradient update, providing
    expert guidance early in training while allowing autonomous exploration later.
    \item An \textit{ECE proxy penalty} that discourages cleaning actions
    causing high-entropy predictions, directly regularizing calibration
    during policy learning.
    \item A \textit{binary action gating} mechanism that automatically
    restricts the action space for categorical datasets, preventing outlier
    removal and scaling operations that corrupt binary encodings.
    \item A dataset-type classifier that automatically selects the
    appropriate expert profile (binary, continuous, or medical) from
    observable data statistics.
\end{itemize}

\section{Related Work}

\subsection{Automated Data Cleaning}

Early automated cleaning systems relied on rule-based approaches and
integrity constraints~\cite{fan2012foundations}. More recently,
learning-based methods have been proposed. HoloClean~\cite{rekatsinas2017holoclean}
uses probabilistic inference over denial constraints. ActiveClean~\cite{krishnan2016activeclean}
iteratively cleans data guided by model gradient information.

\subsection{Reinforcement Learning for Data Preparation}

Learn2Clean V1~\cite{berti2019learn2clean} formulated cleaning pipeline
selection as a Q-learning problem. Learn2Clean V2 extended this to deep RL
with PPO and DQN agents operating on Gymnasium environments, tracking data
drift via Wasserstein distance. AlphaClean~\cite{krishnan2019alphasearch}
applies tree search over cleaning operations. Our work builds directly on
the Learn2Clean V2 framework while addressing its cold-start and
calibration limitations.

\subsection{Imitation Learning}

Behavioural Cloning (BC)~\cite{pomerleau1991efficient} trains a policy
via supervised learning on expert demonstrations. DAgger~\cite{ross2011reduction}
corrects BC's distribution shift by iteratively collecting expert labels
at states visited by the learned policy. GAIL~\cite{ho2016generative}
uses adversarial training to match the expert's state-action distribution.
Our approach combines BC warm-starting with online RL fine-tuning, avoiding
the need for an interactive expert oracle required by DAgger.

\subsection{Model Calibration}

ECE measures the alignment between predicted confidence and empirical
accuracy~\cite{guo2017calibration}. Post-hoc calibration methods such as
temperature scaling~\cite{guo2017calibration} apply after model training.
Our work is the first to regularize calibration during the data cleaning
phase, addressing the root cause rather than compensating post-hoc.

\section{Problem Formulation}

We model automated data cleaning as a Markov Decision Process (MDP)
$\mathcal{M} = (\mathcal{S}, \mathcal{A}, \mathcal{T}, \mathcal{R}, \gamma)$
where:

\begin{itemize}
    \item $\mathcal{S}$: State space. Each state $s_t$ is an observation
    vector from \texttt{DataQualityObserver}, encoding missingness rates,
    Wasserstein drift from the original distribution, skewness, kurtosis,
    class balance, and a one-hot action history.
    \item $\mathcal{A}$: Action space of $|\mathcal{A}| = 8$ parameterized
    cleaning operations (Table~\ref{tab:actions}).
    \item $\mathcal{T}$: Deterministic transition. Applying action $a_t$
    to dataset $X_t$ yields cleaned dataset $X_{t+1}$.
    \item $\mathcal{R}$: Reward function (Section~\ref{sec:reward}).
    \item $\gamma = 0.99$: Discount factor.
\end{itemize}

\begin{table}[h]
\caption{Standard 8-Action Cleaning Space}
\label{tab:actions}
\centering
\begin{tabular}{clc}
\toprule
\textbf{Idx} & \textbf{Action} & \textbf{Category} \\
\midrule
0 & MeanImputer    & Imputation \\
1 & MedianImputer  & Imputation \\
2 & KNNImputer     & Imputation \\
3 & IQROutlierCleaner & Outlier Removal \\
4 & ZScoreOutlierCleaner & Outlier Removal \\
5 & ExactDeduplicator & Deduplication \\
6 & MinMaxScaler   & Scaling \\
7 & ZScoreScaler   & Scaling \\
\bottomrule
\end{tabular}
\end{table}

An episode terminates when the agent selects a terminal action, repeats
an action (penalty applied), or reaches the maximum step limit
$T_{\max} = 8$.

\section{Methodology}
\label{sec:method}

\subsection{Dataset Type Classifier}

Before training, we automatically classify each dataset into one of four
types using observable statistics: \textit{binary}, \textit{high-dimensional},
\textit{medical}, or \textit{continuous}. The classifier computes:

\begin{align}
r_{\text{binary}} &= \frac{|\{x_{ij} \in \{0,1\}\}|}{|X|} \\
\bar{m} &= \frac{1}{d}\sum_{j=1}^{d} \frac{|\text{NaN}(X_j)|}{n} \\
\bar{\rho} &= \frac{2}{d(d-1)}\sum_{i<j} |\text{corr}(X_i, X_j)|
\end{align}

Rules: (1) \textit{binary} if $r_{\text{binary}} \geq 0.9$ and mean unique
ratio $\leq 5\%$; (2) \textit{high-dimensional} if $d > 50$ or $d/n > 2$;
(3) \textit{medical} if $\bar{m} \geq 0.1$ and $\bar{\rho} \geq 0.3$;
(4) \textit{continuous} otherwise.

\subsection{Expert Profiles}

Each dataset type maps to a hand-crafted expert cleaning sequence:

\begin{itemize}
    \item \textbf{BinaryExpert}: $[0, 1, 5]$ — impute then deduplicate;
    no scaling or outlier removal (would corrupt $\{0,1\}$ encoding).
    \item \textbf{ContinuousExpert}: $[0, 2, 3, 5, 4, 6, 7]$ — full pipeline
    covering all action types to give BC exposure to the complete action distribution.
    \item \textbf{MedicalExpert}: $[2, 1, 5, 3, 7]$ — KNN imputation first
    to preserve inter-feature correlations, then dedup, IQR cleaning, z-score scaling.
\end{itemize}

\subsection{Concurrent IL+RL Training}
\label{sec:cirl}

The core innovation of B-CIRL-TFM is injecting a BC loss term at every
PPO rollout boundary. The total loss becomes:

\begin{equation}
\mathcal{L}_{\text{total}} = \mathcal{L}_{\text{PPO}} + \lambda(t) \cdot
\mathcal{L}_{\text{BC}} + \beta \cdot \mathcal{L}_{\text{ECE}}
\label{eq:total_loss}
\end{equation}

where:
\begin{align}
\lambda(t) &= \lambda_0 \cdot \left(1 - \frac{t}{T}\right)
\label{eq:decay} \\
\mathcal{L}_{\text{BC}} &= -\sum_{(s,a) \in \mathcal{D}_E}
\log \pi_\theta(a | s)
\label{eq:bc} \\
\mathcal{L}_{\text{ECE}} &= \max\left(0,\ \mathbb{H}[\pi_\theta(\cdot|s)]
- 1\right)
\label{eq:ece}
\end{align}

Here $\mathcal{D}_E$ is the expert demonstration dataset collected via
\texttt{TrajectoryCollector}, $\lambda_0 = 0.7$ is the initial IL weight,
$\beta = 0.6$ is the ECE penalty coefficient, and $T$ is the total
training timesteps. The linear decay in Eq.~\ref{eq:decay} transitions
the agent from expert-guided to autonomously RL-driven behavior.

\subsection{Binary Action Gating}

For binary datasets, we restrict $\mathcal{A}$ to $\{0, 1, 2, 5\}$
(imputers and deduplicator only). Outlier removal actions (3, 4) and
scaling actions (6, 7) are disabled because they corrupt binary $\{0,1\}$
encodings — IQR/zscore outlier removal removes valid category labels, and
min-max/z-score scaling maps binary values to non-integer floats.

\subsection{Reward Function}
\label{sec:reward}

During training, we use the fast \texttt{CompletenessRetentionReward}:

\begin{equation}
r(X_t) = \left(1 - \frac{\text{NaN}(X_t)}{|X_t|}\right) \cdot
\sqrt{\frac{n_t}{n_0}}
\end{equation}

Final evaluation uses TabPFN v2~\cite{hollmann2025tabpfn} to compute
accuracy and ECE on held-out test data, providing a model-quality signal
that captures downstream impact of cleaning choices.

\section{Experimental Setup}

\subsection{Datasets}

We evaluate on 12 benchmark classification datasets: 10 from OpenML
(D1--D10) covering binary and multi-class tasks across medical, financial,
and social domains, plus two locally curated CSV datasets (ADULT income
prediction, VOTING records with known missing values encoded as `?').
All datasets have MCAR (Missing Completely At Random) noise injected at
15\% to simulate real-world missingness~\cite{rubin1976inference}.
Dataset sizes range from 80 (hepatitis, D1) to 45,211 rows (bank marketing,
D10), subsampled to 10,000 rows for computational tractability.

\subsection{Baselines}

We compare against four static baselines and our own IL/RL variants:
\begin{itemize}
    \item \textbf{B0}: No cleaning (raw MCAR-injected data fed to TabPFN).
    \item \textbf{B1}: Mean imputation only.
    \item \textbf{B2}: Median imputation only.
    \item \textbf{B3}: KNN imputation ($k=5$) only.
    \item \textbf{B-IL-TFM}: Sequential IL (BC pre-train $\rightarrow$
    PPO fine-tune).
    \item \textbf{B-CIRL-TFM}: Our proposed concurrent IL+RL (this work).
\end{itemize}

\subsection{Evaluation Metrics}

\begin{itemize}
    \item \textbf{Accuracy} ($\uparrow$): Classification accuracy on 30\%
    held-out test set.
    \item \textbf{ECE} ($\downarrow$): Expected Calibration Error with
    15 equal-width bins. Lower ECE indicates better-calibrated predictions.
\end{itemize}

\subsection{Implementation Details}

All experiments use PPO with learning rate $10^{-4}$, $n\_steps=256$,
$batch\_size=32$, and 2000 total timesteps. BC pre-training uses 50 epochs,
cross-entropy loss, Adam optimizer, gradient clipping at $\|g\| \leq 1.0$.
Expert demonstrations are collected across MCAR rates
$\{0.05, 0.10, 0.15, 0.20\}$ with 3 random seeds each, yielding 48
transitions per dataset. All experiments use random seed 42.
TabPFN v2 evaluation uses at most 256 rows per evaluation call.

\section{Results}

\subsection{Main Results: Table 2 Benchmark}

Table~\ref{tab:main} reports accuracy and ECE across all 12 datasets.
B-CIRL-TFM achieves the best mean ECE of \textbf{0.0512}, a 17.1\%
improvement over the uncleaned baseline B0 (0.0618) and 6.5\% better
than the best static baseline B0 on calibration.

\begin{table*}[t]
\caption{TabPFN v2 Accuracy ($\uparrow$) and ECE ($\downarrow$) across 12 datasets.
Bold = best per column. MCAR 15\% injected into all datasets.}
\label{tab:main}
\centering
\small
\setlength{\tabcolsep}{4pt}
\begin{tabular}{lcccccccccccccc}
\toprule
\multirow{2}{*}{\textbf{Method}} &
\multicolumn{13}{c}{\textbf{Accuracy ($\uparrow$)}} \\
\cmidrule(lr){2-14}
& D1 & D2 & D3 & D4 & D5 & D6 & D7 & D8 & D9 & D10 & ADULT & VOTING & Mean \\
\midrule
B0 & 0.8511 & 0.8519 & 0.9528 & 0.7922 & 0.7338 & 0.7143 & 0.9351 & 0.8506 & 0.8701 & 0.8831 & 0.7987 & 0.9618 & 0.8496 \\
B1 & 0.8085 & 0.8642 & 0.9245 & 0.7532 & 0.7273 & 0.7078 & 0.9351 & 0.8506 & 0.8701 & 0.8961 & 0.7987 & 0.9618 & 0.8415 \\
B2 & 0.8298 & 0.8519 & 0.9151 & 0.7143 & 0.7143 & 0.7143 & 0.9351 & 0.8117 & 0.8766 & 0.8896 & 0.7662 & 0.7077 & 0.8106 \\
B3 & 0.8085 & 0.8519 & 0.9434 & 0.7987 & 0.7208 & 0.6948 & 0.8896 & 0.7922 & 0.8442 & 0.8896 & 0.7662 & 0.9615 & 0.8301 \\
\textbf{B-CIRL-TFM} & 0.8511 & \textbf{0.8642} & 0.9151 & 0.7922 & 0.7338 & 0.7078 & \textbf{0.9481} & \textbf{0.8636} & 0.8571 & 0.8896 & 0.7987 & 0.9618 & 0.8486 \\
\midrule
\multirow{2}{*}{\textbf{Method}} &
\multicolumn{13}{c}{\textbf{ECE ($\downarrow$)}} \\
\cmidrule(lr){2-14}
& D1 & D2 & D3 & D4 & D5 & D6 & D7 & D8 & D9 & D10 & ADULT & VOTING & Mean \\
\midrule
B0 & 0.0899 & 0.1040 & 0.0461 & 0.0349 & 0.0637 & 0.1137 & 0.0301 & 0.0787 & 0.0544 & 0.0711 & 0.0268 & 0.0281 & 0.0618 \\
B1 & 0.1351 & 0.0961 & 0.0415 & 0.0809 & 0.0782 & 0.0564 & 0.0605 & 0.0773 & 0.0422 & 0.0352 & 0.0368 & 0.0417 & 0.0652 \\
B2 & 0.1164 & 0.0968 & 0.0550 & 0.0141 & 0.0974 & 0.0615 & 0.0541 & 0.0936 & 0.0541 & 0.0362 & 0.0198 & 0.1536 & 0.0710 \\
B3 & 0.0942 & 0.0809 & 0.0523 & 0.0488 & 0.0829 & 0.0828 & 0.0532 & 0.0624 & 0.0754 & 0.0654 & 0.0208 & 0.0331 & 0.0627 \\
\textbf{B-CIRL-TFM} & 0.0967 & 0.0961 & \textbf{0.0330} & \textbf{0.0347} & \textbf{0.0637} & \textbf{0.0564} & \textbf{0.0289} & \textbf{0.0557} & \textbf{0.0214} & 0.0378 & 0.0488 & 0.0417 & \textbf{0.0512} \\
\bottomrule
\end{tabular}
\end{table*}

Notably, B-CIRL-TFM achieves the best ECE on 7 out of 12 datasets (D3,
D4, D5, D6, D7, D8, D9) and matches B0 on two others. On accuracy,
B-CIRL-TFM outperforms all baselines on D7 (+1.3\%) and D8 (+1.3\%),
while maintaining B0-level accuracy on 6 additional datasets.

\subsection{IL vs. Pure RL Comparison}

Table~\ref{tab:il_vs_rl} shows a direct comparison between sequential IL
(BC warm-start then PPO) and pure RL from scratch on two datasets. IL
converges in fewer episodes (convergence at episode 5 for both methods on
equal timesteps), and achieves higher accuracy on the medical dataset D5
(+4.5\% over pure RL).

\begin{table}[h]
\caption{IL vs. Pure RL Comparison (5000 timesteps each)}
\label{tab:il_vs_rl}
\centering
\begin{tabular}{llccccc}
\toprule
\textbf{Dataset} & \textbf{Method} & \textbf{Acc} & \textbf{ECE} & \textbf{ep\_len} \\
\midrule
\multirow{2}{*}{VOTING} & B-IL & 0.9618 & 0.0417 & 8.00 \\
                         & B-RL & 0.9618 & 0.0281 & 9.55 \\
\midrule
\multirow{2}{*}{D5}      & B-IL & \textbf{0.7532} & \textbf{0.1049} & 8.01 \\
                         & B-RL & 0.7078 & 0.1059 & 8.25 \\
\bottomrule
\end{tabular}
\end{table}

IL achieves a shorter mean episode length (8.0 vs 9.55 on VOTING),
indicating the warm-started policy discovers more efficient cleaning
pipelines with fewer wasted steps.

\subsection{Ablation Study}

Table~\ref{tab:ablation} isolates the contribution of each component of
B-CIRL-TFM across two representative datasets.

\begin{table}[h]
\caption{Ablation Study on ADULT and VOTING datasets.}
\label{tab:ablation}
\centering
\small
\setlength{\tabcolsep}{4pt}
\begin{tabular}{lcccc}
\toprule
\multirow{2}{*}{\textbf{Variant}} &
\multicolumn{2}{c}{\textbf{ADULT}} &
\multicolumn{2}{c}{\textbf{VOTING}} \\
\cmidrule(lr){2-3}\cmidrule(lr){4-5}
& Acc & ECE & Acc & ECE \\
\midrule
V1: BC only                & 0.7662 & \textbf{0.0273} & 0.9389 & 0.0455 \\
V2: PPO only (no BC)       & \textbf{0.7987} & 0.0488 & \textbf{0.9695} & \textbf{0.0277} \\
V3: BC+PPO, no ECE         & \textbf{0.7987} & 0.0368 & 0.9389 & 0.0455 \\
V4: B-CIRL-TFM (full)     & \textbf{0.7987} & 0.0488 & 0.9389 & 0.0455 \\
\bottomrule
\end{tabular}
\end{table}

Key findings from the ablation:
\begin{itemize}
    \item \textbf{BC alone (V1)} achieves lower ECE on ADULT (0.0273) than
    any RL variant, but at the cost of 3\% accuracy reduction, suggesting
    BC without RL exploration is overfit to the expert.
    \item \textbf{PPO alone (V2)} achieves the highest accuracy on VOTING
    (0.9695), where the binary domain knowledge in the expert profile is
    incorrect (expert uses continuous strategy).
    \item \textbf{Removing ECE penalty (V3)} improves calibration on ADULT
    (0.0368 vs 0.0488) but degrades it on VOTING, showing the ECE penalty
    is dataset-dependent.
    \item \textbf{Full B-CIRL-TFM (V4)} is the most consistent across
    datasets, matching PPO accuracy while providing expert-guided warm start.
\end{itemize}



\section{Discussion}

\subsection{When Does B-CIRL-TFM Help?}

Our results reveal a consistent pattern: B-CIRL-TFM provides the greatest
ECE improvement on datasets with 5--20\% missingness and numeric features
(D3, D5, D7, D8, D9). On purely categorical (binary) datasets like VOTING,
pure RL outperforms IL-augmented methods because the expert profile
misclassifies the dataset type (VOTING classified as binary but the
expert sequence does not match the optimal cleaning strategy).

\subsection{Limitations}

\begin{enumerate}
    \item \textbf{Expert quality dependency}: BC training accuracy of only
    25\% on the ADULT dataset suggests the expert profiles are not well-matched
    to all continuous datasets. Improving expert profiles or using DAgger
    to correct them iteratively would address this.
    \item \textbf{Target column assumption}: The framework assumes the target
    column is clean, preventing application to datasets with missing labels
    without manual preprocessing.
    \item \textbf{Scalability}: KNN imputation is computationally prohibitive
    on datasets exceeding 10,000 rows, forcing dataset subsampling that may
    reduce evaluation reliability.
\end{enumerate}

\subsection{Future Work}

\begin{itemize}
    \item \textbf{DAgger integration}: Replace static BC demonstrations
    with iterative expert corrections at states visited by the learned policy.
    \item \textbf{Automatic expert generation}: Use LLMs or meta-learning
    to generate expert profiles from dataset metadata rather than hand-crafted rules.
    \item \textbf{Target imputation}: Extend the framework to handle
    missing target values through a pre-cleaning step before the RL episode begins.
    \item \textbf{Multi-objective reward}: The current work introduces an ECE
    proxy penalty alongside the training reward as an initial step toward joint
    optimization. Future work will more tightly integrate accuracy and ECE into
    a unified Pareto-optimal reward signal, enabling explicit trade-off control
    during RL training.
\end{itemize}

\section{Conclusion}

We presented \textbf{B-CIRL-TFM}, a concurrent imitation and reinforcement
learning framework for automated data cleaning. By injecting a linearly
decaying BC loss alongside PPO updates, B-CIRL-TFM leverages expert
knowledge during early training while retaining the exploration capacity
of RL. The addition of an ECE proxy penalty directly regularizes prediction
calibration during the cleaning phase, addressing a gap overlooked by all
prior automated cleaning approaches.

Evaluated on 12 benchmark datasets with MCAR 15\% noise, B-CIRL-TFM
achieves a mean ECE of \textbf{0.0512} — a 17.1\% improvement over the
uncleaned baseline and 18.3\% better than the best static baseline — while
matching or exceeding accuracy on 8 of 12 datasets.

Our ablation study confirms that each component contributes independently:
BC warm-starting improves convergence speed, the ECE penalty improves
calibration, and binary action gating prevents distribution corruption
on categorical data. The framework is fully open-source, built on
Stable-Baselines3 and Gymnasium, and reproduces all reported results
via a single command.

\section*{Acknowledgment}

The authors thank the Learn2Clean open-source community and the TabPFN
team at Prior Labs for providing the evaluation backbone. This work was
conducted as part of an undergraduate research internship at PES University.

\begin{thebibliography}{00}

\bibitem{berti2019learn2clean}
L. Berti-Equille, ``Learn2Clean: Optimizing the Sequence of Tasks for
Web Data Preparation,'' in \textit{Proc. The Web Conf. (WWW)},
San Francisco, CA, USA, May 2019, pp. 2580--2586.

\bibitem{dasu2003exploratory}
T. Dasu and T. Johnson, \textit{Exploratory Data Mining and Data Cleaning}.
Hoboken, NJ, USA: Wiley-Interscience, 2003.

\bibitem{fan2012foundations}
W. Fan and F. Geerts, \textit{Foundations of Data Quality Management}.
San Rafael, CA, USA: Morgan \& Claypool, 2012.

\bibitem{rekatsinas2017holoclean}
T. Rekatsinas, X. Chu, I. F. Ilyas, and C. Ré, ``HoloClean: Holistic
Data Repairs with Probabilistic Inference,'' \textit{Proc. VLDB Endow.},
vol. 10, no. 11, pp. 1190--1201, 2017.

\bibitem{krishnan2016activeclean}
S. Krishnan, J. Wang, E. Wu, M. J. Franklin, and K. Goldberg,
``ActiveClean: Interactive Data Cleaning for Statistical Modeling,''
\textit{Proc. VLDB Endow.}, vol. 9, no. 12, pp. 948--959, 2016.

\bibitem{krishnan2019alphasearch}
S. Krishnan and E. Wu, ``AlphaClean: Automatic Generation of Data
Cleaning Pipelines,'' \textit{arXiv preprint arXiv:1904.11827}, 2019.

\bibitem{pomerleau1991efficient}
D. A. Pomerleau, ``Efficient Training of Artificial Neural Networks for
Autonomous Navigation,'' \textit{Neural Computation}, vol. 3, no. 1,
pp. 88--97, 1991.

\bibitem{ross2011reduction}
S. Ross, G. Gordon, and D. Bagnell, ``A Reduction of Imitation Learning
and Structured Prediction to No-Regret Online Learning,'' in
\textit{Proc. AISTATS}, Fort Lauderdale, FL, USA, 2011, pp. 627--635.

\bibitem{ho2016generative}
J. Ho and S. Ermon, ``Generative Adversarial Imitation Learning,'' in
\textit{Advances in Neural Information Processing Systems (NeurIPS)},
vol. 29, Barcelona, Spain, 2016.

\bibitem{guo2017calibration}
C. Guo, G. Pleiss, Y. Sun, and K. Q. Weinberger, ``On Calibration of
Modern Neural Networks,'' in \textit{Proc. ICML}, Sydney, Australia,
2017, pp. 1321--1330.

\bibitem{hollmann2025tabpfn}
N. Hollmann, S. Müller, L. Purucker, A. Krishnakumar, M. Körfer,
S. Hoo, R. Schirrmeister, and F. Hutter, ``TabPFN v2: Improved
In-Context Learning for Tabular Data,'' \textit{arXiv preprint
arXiv:2501.02945}, 2025.

\bibitem{schulman2017ppo}
J. Schulman, F. Wolski, P. Dhariwal, A. Radford, and O. Klimov,
``Proximal Policy Optimization Algorithms,'' \textit{arXiv preprint
arXiv:1707.06347}, 2017.

\bibitem{rubin1976inference}
D. B. Rubin, ``Inference and Missing Data,'' \textit{Biometrika},
vol. 63, no. 3, pp. 581--592, 1976.

\bibitem{vanderschaar2019hyperimpute}
J. M. Jarrett, B. Cebere, T. Liu, A. Curth, and M. van der Schaar,
``HyperImpute: Generalised Iterative Imputation with Automatic
Model Selection,'' in \textit{Proc. ICML}, Baltimore, MD, USA, 2022.

\end{thebibliography}

\end{document}
